\documentclass[10pt, oneside, final]{report}
\usepackage{geometry}
\geometry{a4paper, left=10mm, right=10mm, top=10mm, bottom=20mm}
\usepackage[utf8]{inputenc}
\usepackage{t1enc}
\usepackage[MeX]{polski}
\usepackage{graphicx}
\usepackage{parskip}
\usepackage[table]{xcolor}
\usepackage{booktabs}
\usepackage{ltablex}

% ustawianie pozycji elementow na stronie
\usepackage[absolute, overlay]{textpos}
\setlength{\TPHorizModule}{\textwidth}\setlength{\TPVertModule}{\textheight}
\textblockorigin{10mm}{10mm}

% Czcionka
\usepackage{helvet}
\renewcommand{\familydefault}{\sfdefault}

% Wysokosc pól sprzedawca, nabywca
\newcommand{\heightA}{0.14}

% Numerowanie stron
\usepackage{lastpage}
\usepackage{fancyhdr} % to change header and footers
\usepackage{blindtext} % to quickly get a full document
\pagestyle{fancy} % Turn on the style
\fancyhf{} % Start with clearing everything in the header and footer
% Set the right side of the footer to be the page number
\renewcommand{\headrulewidth}{0pt}
\fancyfoot[R]{\thepage \ / \pageref{LastPage}}

\makeatletter
\setlength{\@fptop}{0pt}
\makeatother

\newcounter{invoicecnt}
\setcounter{invoicecnt}{1}
\newcommand{\invoiceline}[6]{\arabic{invoicecnt} & #1 & #2 & #3 & #4\% & #5 & #6\stepcounter{invoicecnt}\\}
\input{invoice_data.tex}

\begin{document}
% Logo
%\hspace*{-1.8em}
\includegraphics[width=16em]{logo.png}\\

% Tytuł faktury
\begin{textblock}{0.4}(0.5,0.01)
	\large
	\textbf{FAKTURA VAT nr \invoiceNumber}\vspace{0.1em}\\
	\normalsize
	Data wystawienia: \dateInvoice\\
	Data sprzedaży: \dateSold\\
	Sposób zapłaty: przelew\vspace{0.2em}\\
	\textbf{Termin płatności: \datePayment}
\end{textblock}

% Dane sprzedawcy
\begin{textblock}{0.4}(0.0,\heightA)
	Sprzedawca:\vspace{0.2em}\\
	\textbf{ElkaOkna Sp. z o.o.}\\
	ul. Nowowiejska 15/19\\
	00-665 Warszawa\\
	NIP: 525-000-58-34\\
	tel: 22 234 74 97
\end{textblock}

% Dane nabywcy
\begin{textblock}{0.4}(0.5,\heightA)
	Nabywca:\vspace{0.2em}\\
	\textbf{\lineA}\\
	\lineB\\
	\lineC\\
	\lineD\\
	\lineE
\end{textblock}

% Lista zakupów
\vspace{0.11\textheight}
\setlength\extrarowheight{3pt}
\setlength{\tabcolsep}{3pt}
\setlength\aboverulesep{0pt}
\setlength\belowrulesep{0pt}
\keepXColumns
{
\rowcolors{2}{white}{gray!15}
\begin{tabularx}{\linewidth}{lXrrrrr}
	\rowcolor{white}
	\multicolumn{7}{l}{\textbf{Lista zamówionych produktów i usług:}}\\
	\toprule
	\rowcolor{gray!35}
	& &	& Wartość & Stawka & Wartość & Wartość\\
	\rowcolor{gray!35}
	Lp. & Nazwa produktu / usługi & Ilość & NETTO & VAT & VAT & BRUTTO\\
	\midrule[1pt]
	\endhead
	\invoiceLines
	\if\delivery1
	\midrule
	\rowcolor{blue!5}
	\arabic{invoicecnt} & Dostawa & & \deliveryNetto & 8\% & \deliveryVat & \deliveryBrutto\stepcounter{invoicecnt}\\
	\fi
	
	\if\assembly1
	\midrule
	\rowcolor{blue!5}
	\arabic{invoicecnt} & Montaż & & 2 & 8\% & 0,16 & 2,16\stepcounter{invoicecnt}\\
	\fi
	
	\bottomrule
	\rowcolor{white}
	\multicolumn{2}{r}{\textbf{Razem}:} & \sumCounts & \sumNetto & & \sumVat & \sumBrutto\\
\end{tabularx}
}

\begin{figure}[ht!]
	\textbf{Dane do przelewu:}\\
	\vspace{-2em}
	\begin{tabularx}{\linewidth}{lX}
		Konto: & BANK POCZTOWY SA\\
		& 09 1120 2123 2100 0000 3451 0007\\
		Do zapłaty: & \sumBrutto \ zł\\
		Słownie: & \sumTextZl \ zł. \sumTextGr \ gr.\\ \\
		Wystawił: & Kierownik Biura Sprzedaży
	\end{tabularx}
\end{figure}


\end{document}
